\documentclass[11pt]{article}
\usepackage[margin=1in]{geometry}
\usepackage{enumitem}
% =============================================================================
% Preambles/header.tex — shared LaTeX/Beamer preamble for this project
%
% Use from a Beamer lecture by prepending:
%     \input{header}
% and compiling with TEXINPUTS=../Preambles:$TEXINPUTS (the /compile-latex
% skill does this automatically).
%
% PALETTE CONTRACT. Color names defined here MUST match the names in
% Quarto/theme-template.scss so Beamer and Quarto renderings use the same
% palette. The scripts/check-palette-sync.sh script enforces this.
%
% To customize: edit the HEX values below AND the matching $variable in
% Quarto/theme-template.scss, then run scripts/check-palette-sync.sh.
% =============================================================================

% -----------------------------------------------------------------------------
% Engine detection + encoding
%
% Our /compile-latex skill uses XeLaTeX, where inputenc is unnecessary and
% can produce encoding warnings. Only load inputenc under pdfTeX.
% -----------------------------------------------------------------------------
\usepackage{iftex}
\ifPDFTeX
  \usepackage[utf8]{inputenc}
\fi

\usepackage{amsmath, amssymb, amsthm}
\usepackage{booktabs}
\usepackage{graphicx}

% -----------------------------------------------------------------------------
% PALETTE — mirrors Quarto/theme-template.scss variable names.
% Load xcolor and define colors BEFORE anything that references them
% (notably hyperref, hypersetup, and the Beamer color assignments).
% -----------------------------------------------------------------------------
\usepackage{xcolor}

\definecolor{primary-blue}{HTML}{012169}      % headings, accents
\definecolor{primary-gold}{HTML}{B9975B}      % emphasis, borders
\definecolor{highlight-yellow}{HTML}{F2A900}  % markers, alerts
\definecolor{light-bg}{HTML}{E8EDF5}          % title / section backgrounds
\definecolor{jet}{HTML}{1A1A1A}               % body text

% Semantic colors (used in diagrams and annotations)
\definecolor{positive}{HTML}{15803D}          % good / observed / identified
\definecolor{negative}{HTML}{B91C1C}          % bad / problematic / bias
\definecolor{neutral}{HTML}{525252}           % reference / context

% Highlight accents (match .hi-* classes in the SCSS)
\definecolor{hi-slate}{HTML}{314F4F}
\definecolor{hi-green}{HTML}{2E8B57}
\definecolor{hi-red}{HTML}{C41E3A}

% Semantic aliases so skill templates and snippets can write
% \textcolor{accent}{...} without hard-coding "primary-blue".
\colorlet{accent}{primary-blue}
\colorlet{accent2}{primary-gold}

% -----------------------------------------------------------------------------
% Hyperref — loaded AFTER xcolor + palette so link colors resolve.
% -----------------------------------------------------------------------------
\usepackage{hyperref}
\hypersetup{colorlinks=true, linkcolor=primary-blue, urlcolor=primary-blue,
            citecolor=primary-blue, pdfborder={0 0 0}}

% -----------------------------------------------------------------------------
% Course code — set once per deck (after \input{header}, before \begin{document})
% via \coursecode{CS401}. Shown in the footer on every slide so a deck's course
% is identifiable without opening the title page. Empty by default.
% -----------------------------------------------------------------------------
\makeatletter
\newcommand{\coursecode}[1]{\gdef\@coursecodetext{#1}}
\gdef\@coursecodetext{}
\makeatother

% -----------------------------------------------------------------------------
% Beamer theme assignments (only apply under Beamer).
%
% We detect Beamer via \@ifclassloaded{beamer}, which is the documented,
% non-internal way. This must be wrapped in \makeatletter/\makeatother
% because \@ifclassloaded contains an @-catcode character.
% -----------------------------------------------------------------------------
\makeatletter
\@ifclassloaded{beamer}{%
  \usefonttheme{professionalfonts}
  \setbeamercolor{structure}{fg=primary-blue}
  \setbeamercolor{frametitle}{fg=primary-blue}
  \setbeamercolor{title}{fg=primary-blue}
  \setbeamercolor{subtitle}{fg=primary-gold}
  \setbeamercolor{author}{fg=primary-blue}
  \setbeamercolor{institute}{fg=neutral}
  \setbeamercolor{date}{fg=primary-gold}
  \setbeamercolor{itemize item}{fg=primary-blue}
  \setbeamercolor{itemize subitem}{fg=primary-gold}
  \setbeamercolor{alerted text}{fg=primary-gold}
  \setbeamercolor{block title}{fg=white, bg=primary-blue}
  \setbeamercolor{block body}{bg=light-bg}
  % Semantic block variants: block=definition/concept (blue), exampleblock=
  % worked example (green/positive), alertblock=key takeaway or caution (gold).
  % Keep to <=2 colored boxes per slide across ALL three variants (INV-7).
  \setbeamercolor{block title example}{fg=white, bg=positive}
  \setbeamercolor{block body example}{bg=light-bg}
  \setbeamercolor{block title alerted}{fg=white, bg=primary-gold}
  \setbeamercolor{block body alerted}{bg=light-bg}

  % Minimal footer: course code (left) + page X / N (right), no navigation clutter
  \setbeamertemplate{navigation symbols}{}
  \setbeamertemplate{footline}{%
    \color{neutral}\footnotesize\hspace{0.7em}\@coursecodetext\hfill
    \insertframenumber/\inserttotalframenumber\hspace{0.7em}\vspace{0.3em}}
}{}
\makeatother

% -----------------------------------------------------------------------------
% TikZ setup — libraries the snippet gallery uses
% -----------------------------------------------------------------------------
\usepackage{tikz}
\usetikzlibrary{arrows.meta, positioning, calc, decorations.pathreplacing,
                fit, shapes.geometric, backgrounds}

% Shared TikZ styles — snippets and hand-written diagrams can pick these up
% via `\begin{tikzpicture}[every node/.style={...}]` or by naming specific
% styles (e.g., `\node[dag-node] (X) at (0,0) {X};`).
\tikzset{
  % Node styles for causal DAGs and similar
  dag-node/.style = {
    draw=primary-blue, fill=light-bg, rounded corners=2pt,
    minimum width=1.2cm, minimum height=0.9cm, align=center, font=\small
  },
  decision-node/.style = {
    draw=primary-gold, fill=white, diamond, aspect=1.8,
    minimum width=1.8cm, minimum height=1.2cm, align=center, font=\small,
    inner sep=1pt
  },
  % Edge styles — observed vs counterfactual
  observed-edge/.style = {-{Latex[length=2.5mm]}, thick, draw=jet},
  counterfactual-edge/.style = {-{Latex[length=2.5mm]}, thick, draw=neutral, dashed},
  confound-edge/.style = {-{Latex[length=2.5mm]}, thick, draw=negative, dashed},
  % Data-point styles for plots
  observed-dot/.style = {fill=primary-blue, draw=primary-blue, circle, inner sep=1.2pt},
  counterfactual-dot/.style = {fill=white, draw=neutral, circle, inner sep=1.2pt}
}

% -----------------------------------------------------------------------------
% Convenience macros
% -----------------------------------------------------------------------------
\newcommand{\muted}[1]{\textcolor{neutral}{#1}}
\newcommand{\key}[1]{\textcolor{primary-gold}{\textbf{#1}}}
\newcommand{\good}[1]{\textcolor{positive}{#1}}
\newcommand{\bad}[1]{\textcolor{negative}{#1}}

% Standout frame macro: full-bleed dark background for section transitions.
% Beamer-only (uses \begin{frame}/\setbeamercolor/\usebeamerfont) -- do not
% call from an article-class file (e.g. Notes/<CODE>/*-notes.tex). \muted,
% \key, \good, \bad, and \coursecode above are plain xcolor/gdef and are
% safe to use from either class; verified when Notes/article-class support
% was added.
% Expand: \transitionslide{Your transition title}
\newcommand{\transitionslide}[1]{%
  \begingroup
  \setbeamercolor{background canvas}{bg=primary-blue}
  \begin{frame}[plain]
    \centering
    \vfill
    {\usebeamerfont{title}\color{white}\Large #1\par}
    \vfill
  \end{frame}
  \endgroup
}

% Section divider frame (Beamer-only), an alternative to \transitionslide with a
% darker two-line style: near-black (jet) background, a small white label line
% above a larger gold title, and a thin gold rule along the bottom edge.
% Expand: \sectiondivider{Section 2}{Pointers}
\newcommand{\sectiondivider}[2]{%
  \begingroup
  \setbeamercolor{background canvas}{bg=jet}
  \begin{frame}[plain]
    \vspace*{3.0cm}
    {\color{white}\ttfamily\large #1\par}
    \vspace{0.5cm}
    {\color{primary-gold}\LARGE #2\par}
    \begin{tikzpicture}[remember picture, overlay]
      \draw[primary-gold, line width=2.5pt]
        ($(current page.south west)+(0,0.1)$) -- ($(current page.south east)+(0,0.1)$);
    \end{tikzpicture}
  \end{frame}
  \endgroup
}

% -----------------------------------------------------------------------------
% Channel watermark — "Concepts That Click" (YouTube).
%
% Article-class files (CompetitiveExam Q/A sets, Notes, Assignments) get a
% light diagonal draft watermark on every page plus a centered footer naming
% the channel. Beamer decks get a subtle TikZ background watermark instead
% (the footline already carries the course code). Centralized here so every
% MCQ PDF is branded without per-file edits.
% -----------------------------------------------------------------------------
\makeatletter
\@ifclassloaded{article}{%
  \usepackage{draftwatermark}
  \SetWatermarkText{Concepts That Click}
  \SetWatermarkScale{1}
  \SetWatermarkFontSize{2cm}
  \SetWatermarkLightness{0.86}
  \SetWatermarkAngle{45}
  \usepackage{fancyhdr}
  \pagestyle{fancy}
  \fancyhf{}
  \fancyfoot[C]{\footnotesize\textcolor{neutral}{Concepts That Click~|~YouTube: \href{https://www.youtube.com/@concepts.thatclick}{@concepts.thatclick}}}
  \renewcommand{\headrulewidth}{0pt}
  \renewcommand{\footrulewidth}{0pt}
  \fancypagestyle{plain}{%
    \fancyhf{}%
    \fancyfoot[C]{\footnotesize\textcolor{neutral}{Concepts That Click~|~YouTube: \href{https://www.youtube.com/@concepts.thatclick}{@concepts.thatclick}}}%
    \renewcommand{\headrulewidth}{0pt}%
    \renewcommand{\footrulewidth}{0pt}%
  }
}{}
\@ifclassloaded{beamer}{%
  \setbeamertemplate{background}{%
    \begin{tikzpicture}[remember picture, overlay]
      \node[rotate=45, opacity=0.055, align=center]
        at (current page.center)
        {\Huge\textcolor{neutral}{Concepts That Click}};
    \end{tikzpicture}%
  }
}{}
\makeatother 
\coursecode{JKSSB}
\renewcommand{\thesection}{19.\arabic{section}}
\newtheorem{example}{Example}[section]
\newenvironment{definitionbox}[1]{%
  \par\noindent\textbf{Definition (#1).}\ \itshape
}{\par\vspace{0.3em}}
\title{Operating System --- Detailed Lecture Notes}
\author{JKSSB Computer Awareness}
\date{\today}

\begin{document}
\maketitle

\section{What an operating system is}
A computer's hardware cannot be used directly by a person or by ordinary
application programs. Something must sit between the user and the machine,
accept the user's instructions, and control the hardware on the user's behalf.
That something is the \textbf{operating system}. An operating system (OS) is
the main piece of \textbf{system software} that controls the whole computer. It
starts the machine, runs other programs, and manages every resource. It makes
the computer usable, because without an OS the hardware alone cannot run
application software.

\begin{definitionbox}{Operating system}
The main system software that sits between the user and the computer hardware,
provides a user-friendly interface, and manages the computer's resources such
as the processor, memory, storage, and input/output devices.
\end{definitionbox}

\begin{center}
\begin{tabular}{p{0.24\textwidth}p{0.66\textwidth}}
\toprule
\textbf{Term} & \textbf{Meaning in this lesson} \\
\midrule
Kernel & The core part of the OS that interacts directly with the hardware. \\
Booting & Starting the computer and loading the OS into main memory. \\
GUI & Graphical User Interface --- interaction through icons, windows, and clicks. \\
CLI & Command Line Interface --- interaction by typing text commands. \\
Process & A program that is currently executing. \\
Multitasking & Running more than one program at a time. \\
Time-sharing & Sharing one computer among several users at the same time. \\
\bottomrule
\end{tabular}
\end{center}

The rest of this lesson uses these terms in one consistent sense. The
\textbf{kernel} is the privileged core of the OS. \textbf{Booting} is the
startup sequence. A \textbf{GUI} is a graphical interface and a \textbf{CLI}
is a text-command interface. A \textbf{process} is a running program.
\textbf{Multitasking} and \textbf{time-sharing} are two ways of sharing the
CPU.

\section{The OS as interface and resource manager}
The operating system has \textbf{two main jobs}. The first is to provide a
\textbf{user-friendly interface}: the user gives instructions and receives
results through the OS, without needing to understand the hardware. This is
the definition of the OS as an interface layer, and it is a regularly tested
idea (PYQ-21). The second job is to act as a \textbf{resource manager}: the OS
shares the CPU, main memory, storage, and devices among all the programs that
are running, so that each one gets a fair turn and none interferes with the
others.

It is important to know \textbf{where} the operating system is. The OS is
stored permanently on secondary storage, such as an SSD or HDD. At startup it
is loaded into \textbf{primary storage}, that is, main memory. While the
computer is running, the OS is \emph{resident in primary storage} (PYQ-22): it
stays in main memory for the whole session so that it can respond instantly
whenever the hardware or an application needs it.

\begin{example}
An item states that the OS is kept in secondary storage while the computer is
running. This is wrong. The OS is stored on secondary storage when the machine
is off, but once the computer starts it is copied into primary storage (main
memory) and remains there while running. The wrong option confuses the place
where the OS is stored with the place where it runs.
\end{example}

\section{Functions of an operating system}
The work of an OS is usually grouped into a few large functions, each of which
is a favourite exam topic:

\begin{enumerate}
  \item \textbf{Process management}: creating processes, scheduling the CPU
    among them, and removing them when they finish.
  \item \textbf{Memory management}: deciding how much main memory each
    program gets and keeping programs apart.
  \item \textbf{File management}: organising data into files and folders and
    handling their creation, reading, writing, and deletion.
  \item \textbf{Device management}: controlling input and output devices
    through device drivers.
  \item \textbf{Security}: managing user accounts, logins, and permissions, and
    applying updates.
  \item \textbf{User interface}: presenting either a GUI or a CLI to the user.
\end{enumerate}

\section{The kernel}
The \textbf{kernel} is the \textbf{core} of the operating system. It is the
part of the OS that interacts \textbf{directly with the hardware}, and it
always stays in main memory while the computer runs. The kernel provides the
most important services: process scheduling, memory handling, and device
control. Other software, including the user interface, relies on the kernel to
reach the hardware.

\begin{definitionbox}{Kernel}
The central, privileged part of the operating system that manages the hardware
and provides the fundamental services on which all other software depends.
\end{definitionbox}

\section{User interfaces: GUI and CLI}
Two broad kinds of interface are used to work with a computer. A
\textbf{Graphical User Interface (GUI)} lets the user interact through
pictures, icons, windows, menus, and clicks; it is easier for beginners. A
\textbf{Command Line Interface (CLI)} requires the user to type text commands;
it gives precise control but demands that the commands be known. Windows,
macOS, and Linux desktops are examples of GUIs, while the command prompt and
the terminal are examples of CLIs.

\begin{center}
\begin{tabular}{p{0.24\textwidth}p{0.30\textwidth}p{0.30\textwidth}}
\toprule
\textbf{Point} & \textbf{GUI} & \textbf{CLI} \\
\midrule
Full form & Graphical User Interface & Command Line Interface \\
Input method & Clicks, icons, windows & Typed text commands \\
Ease of use & Easier for beginners & Needs the commands to be known \\
Example & Windows desktop, icons & Command prompt, terminal \\
\bottomrule
\end{tabular}
\end{center}

\section{Booting the computer}
\textbf{Booting} is the process of starting a computer and loading the
operating system into main memory. A \textbf{cold boot} starts from the
power-off state, while a \textbf{warm boot} restarts a computer that is already
on. The startup firmware, \textbf{BIOS} (Basic Input/Output System) or its
modern replacement \textbf{UEFI}, first checks the hardware and then finds and
loads the operating system from secondary storage. BIOS/UEFI initialises the
hardware and hands control to the OS; it is not the kernel and it is not a
device driver.

\section{Process management}
A \textbf{process} is a program that is currently executing. A program sitting
on disk is passive; once it is loaded and running it becomes an active
process. The operating system creates processes, gives each a turn on the CPU
through \textbf{CPU scheduling}, lets them communicate where necessary, and
removes them when they finish. Scheduling decides which process runs next and
for how long, which is what allows several programs to make progress.

\section{Memory management}
The operating system decides \textbf{how much main memory} each program
receives. It keeps one program's memory \textbf{separate} from another's so
that programs cannot overwrite each other's data. It frees memory when a
program ends so that the space can be reused. Where needed, the OS uses
\textbf{virtual memory} to give programs more space than the physical RAM
provides.

\section{File management}
The OS organises data on storage devices into \textbf{files} and
\textbf{folders}. It handles creating, reading, writing, renaming, and deleting
files, and it keeps a \textbf{directory} that records where each file is
stored. It also controls \textbf{permissions}: which accounts may open, change,
or delete a given file. Together these tasks are called the \textbf{file
system}, the way the OS names, stores, and finds files.

\section{Device management}
The operating system controls all \textbf{input/output devices}: the keyboard,
mouse, printer, disk, and so on. Each device communicates with the OS through a
\textbf{device driver}. A device driver is \textbf{software}, not hardware ---
a point that is regularly tested (TRAP-10). The OS also uses \textbf{spooling}
to queue print jobs and feed the printer at the printer's own speed, so that
the computer does not have to wait for a slow device to finish.

\section{Multitasking and time-sharing}
\textbf{Multitasking} means that the operating system runs \textbf{more than
one program at a time} (PYQ-23). The CPU switches between programs so quickly
that they appear to run together, even though a single CPU core does only one
thing at any instant. \textbf{Time-sharing} lets \textbf{many users} share one
computer at the same time, each receiving a slice of CPU time in turn. Both
multitasking and time-sharing depend on the OS \textbf{scheduler}.

\begin{example}
An item defines multitasking as using more than one monitor, or as storing more
than one file on a disk. Both are wrong. Multitasking is about the
\textbf{processor} running more than one program at a time. The number of
monitors and the number of files stored are unrelated to multitasking.
\end{example}

\section{Security, user accounts and updates}
The operating system provides \textbf{user accounts}, each with a username and
password. \textbf{Authentication} (the login step) verifies who is using the
computer. \textbf{Permissions} decide which files and settings each account may
use, so that one user cannot disturb another's work. Finally, system and
software \textbf{updates} should be applied regularly, because they fix
security holes and add improvements. Keeping the OS up to date is one of the
simplest ways to protect the machine.

\section{Examples of operating systems}
Common operating systems include \textbf{Microsoft Windows}, the open-source
\textbf{Linux} family, and Apple's \textbf{macOS} on desktops and laptops, plus
the mobile operating systems \textbf{Android} and \textbf{iOS}. Android is
built on the Linux kernel. An application such as MS Word or a web browser is
\emph{not} an operating system; it runs on top of one.

\section{Exam retrieval and common confusions}
The following distinctions prevent most errors on this topic.
\begin{enumerate}
  \item The OS is \textbf{system software}, not application software. A device
    driver is \textbf{software}, not hardware (TRAP-10).
  \item While running, the OS is resident in \textbf{primary storage} (main
    memory); it is only \emph{stored} on secondary storage (PYQ-22).
  \item The OS provides a \textbf{user-friendly interface} between the user and
    the hardware (PYQ-21).
  \item \textbf{Multitasking} means running more than one program at a time;
    \textbf{time-sharing} shares one machine among several users (PYQ-23).
  \item \textbf{BIOS/UEFI} initialises the hardware and loads the OS; it is not
    the kernel and not a device driver (TRAP-24).
  \item A \textbf{program} on disk is passive; once running it becomes a
    \textbf{process}.
\end{enumerate}

\begin{example}
An item asks where the operating system is loaded when the computer starts.
Trace the startup sequence: the machine is switched on, BIOS/UEFI checks the
hardware, and the OS is loaded from secondary storage into \textbf{primary
storage} (main memory), where it remains while the computer runs. If the stem
instead asks which component is the core of the OS, the answer is the
\textbf{kernel}.
\end{example}

\subsection*{Exercises}
\begin{enumerate}[label=\arabic*.]
  \item Define an operating system and state its two main roles.
  \item Where is the operating system stored, and where is it resident while
    the computer is running?
  \item List four core functions of an operating system.
  \item What is the kernel, and why is it important?
  \item Distinguish a GUI from a CLI.
  \item What does multitasking mean, and how does it differ from
    time-sharing?
  \item State three ways an operating system helps to keep a computer secure.
\end{enumerate}

\subsection*{Solutions}
\begingroup\small
\begin{enumerate}[label=\arabic*.]
  \item An operating system is the main system software that sits between the
    user and the hardware. Its two main roles are to provide a user-friendly
    interface and to manage the computer's resources.
  \item It is stored permanently on secondary storage (such as an SSD or HDD)
    and is resident in primary storage (main memory) while the computer is
    running.
  \item Any four of: process management, memory management, file management,
    device management, security, and the user interface.
  \item The kernel is the core part of the operating system that interacts
    directly with the hardware. It is important because it provides the
    fundamental services --- scheduling, memory handling, and device control
    --- on which all other software depends.
  \item A GUI (Graphical User Interface) uses icons, windows, and clicks and is
    easier for beginners; a CLI (Command Line Interface) requires the user to
    type text commands and gives precise control.
  \item Multitasking means the OS runs more than one program at a time by
    switching the CPU among them. Time-sharing means many users share one
    computer at the same time, each getting a slice of CPU time.
  \item Three of: user accounts with usernames and passwords; authentication at
    login; permissions that limit access to files and settings; and regular
    security updates.
\end{enumerate}
\endgroup
\end{document}
